%%%%%%%%%%%%%%%%%%%%%%%%%%%%%%%%%%%%%%%%%%%%%%%%%%%%%%%%%
% Presentasi Beamer untuk RPS Optimisasi Teknik Elektro
% Minggu ke-2: Optimisasi Linier (LP) & Metode Grafis
% Program Studi S1 Teknik Elektro - Universitas Bengkulu (2026)
%%%%%%%%%%%%%%%%%%%%%%%%%%%%%%%%%%%%%%%%%%%%%%%%%%%%%%%%%

\documentclass[10pt,aspectratio=169]{beamer}

\usetheme{Madrid}
\usecolortheme{seahorse}
\setbeamertemplate{navigation symbols}{}

\definecolor{unibblue}{RGB}{0, 51, 102}
\definecolor{unibgold}{RGB}{212, 175, 55}
\definecolor{darkgreen}{RGB}{34, 139, 34}

\setbeamercolor{structure}{fg=unibblue}
\setbeamercolor{palette primary}{bg=unibblue,fg=white}
\setbeamercolor{palette secondary}{bg=unibblue!80,fg=white}
\setbeamercolor{palette tertiary}{bg=unibblue!60,fg=white}
\setbeamercolor{titlelike}{parent=palette primary,fg=white}

\usepackage[utf8]{inputenc}
\usepackage[indonesian]{babel}
\usepackage{amsmath, amssymb, amsfonts}
\usepackage{graphicx}
\usepackage{booktabs}
\usepackage{listings}
\usepackage{tcolorbox}
\tcbuselibrary{skins,breakable}

\lstset{
  language=Python,
  basicstyle=\ttfamily\tiny,
  keywordstyle=\color{unibblue}\bfseries,
  commentstyle=\color{darkgreen}\itshape,
  stringstyle=\color{red!70!black},
  numbers=left,
  numberstyle=\tiny\color{gray},
  numbersep=4pt,
  breaklines=true,
  frame=single,
  rulecolor=\color{unibblue!30},
  tabsize=4
}

\newcommand{\R}{\mathbb{R}}

\title[Optimisasi TE: Pertemuan 2]{Optimisasi untuk Teknik Elektro (TKE-41XX)}
\subtitle{Pertemuan 2: Optimisasi Linier (LP) dan Metode Grafis 2D}
\author[N. Daratha]{Ir. Novalio Daratha S.T., M.Sc., Ph.D.}
\institute[Teknik Elektro UNIB]{Program Studi S1 Teknik Elektro\\Fakultas Teknik --- Universitas Bengkulu}
\date{Semester Ganjil 2026}

\begin{document}

\begin{frame}
    \titlepage
\end{frame}

\begin{frame}{Agenda Pembelajaran Minggu ke-2}
    \begin{tcolorbox}[colback=unibblue!5,colframe=unibblue,title=\textbf{Tujuan Pembelajaran (Sub-CPMK 2)}]
        Mahasiswa mampu menyelesaikan masalah Linear Programming (LP) 2D secara grafis dan komputasional, serta menganalisis geometri ruang solusi poligon feasibel.
    \end{tcolorbox}

    \vspace{0.3cm}
    \begin{enumerate}
        \item \textbf{Karakteristik Masalah LP}: Sifat linieritas, proporsionalitas, dan aditivitas.
        \item \textbf{Metode Grafis 2D}: Penggambaran kendala pertidaksamaan dan pembentukan poligon feasibel.
        \item \textbf{Garis Isoprofit / Isocost}: Pencarian titik pojok ekstrim optimal (*Corner Point Feasible*).
        \item \textbf{Studi Kasus Sederhana}: Alokasi beban transformator daya.
        \item \textbf{Implementasi Komputasi}: Penggunaan `scipy.optimize.linprog`.
    \end{enumerate}
\end{frame}

\section{Metode Grafis 2D}

\begin{frame}{Geometri Ruang Solusi Linear Programming}
    \begin{columns}
        \column{0.5\textwidth}
        \begin{block}{Sifat Poligon Feasibel}
            \begin{itemize}
                \item Setiap kendala linier membagi bidang $\R^2$ menjadi dua separuh bidang (*half-space*).
                \item Irisan dari seluruh kendala membentuk **Poligon Konveks Feasibel**.
                \item **Teorema Utama LP**: Solusi optimal selalu terletak pada minimal satu titik pojok ekstrim (*Corner Point Solution*).
            \end{itemize}
        \end{block}

        \column{0.46\textwidth}
        \begin{tcolorbox}[colback=unibgold!10,colframe=unibgold!80!black,title=\textbf{Formulasi Standar LP}]
            \[
                \min_{\mathbf{x}} \mathbf{c}^T \mathbf{x}
            \]
            \[
                \text{s.t.} \quad A \mathbf{x} \le \mathbf{b}, \quad \mathbf{x} \ge 0
            \]
        \end{tcolorbox}
    \end{columns}
\end{frame}

\begin{frame}[fragile]{Implementasi LP dengan scipy.optimize.linprog}
\begin{lstlisting}
import numpy as np
from scipy.optimize import linprog

# Maksimalkan Z = 500 x1 + 800 x2  => Minimalkan -500 x1 - 800 x2
c = [-500, -800]

# Kendala Pertidaksamaan A_ub * x <= b_ub
A_ub = [
    [2, 1],  # 2 x1 + 1 x2 <= 40
    [1, 3]   # 1 x1 + 3 x2 <= 90
]
b_ub = [40, 90]

# Batas variabel x1 >= 0, x2 >= 0
x_bounds = (0, None)

res = linprog(c, A_ub=A_ub, b_ub=b_ub, bounds=[x_bounds, x_bounds], method='highs')

print(f"x1 Optimal = {res.x[0]:.2f}, x2 Optimal = {res.x[1]:.2f}")
print(f"Keuntungan Maksimum = Rp {-res.fun:.2f}")
\end{lstlisting}
\end{frame}

\begin{frame}{Rangkuman Pertemuan 2}
    \begin{tcolorbox}[colback=unibblue!10,colframe=unibblue,title=\textbf{Poin Kunci}]
        \begin{itemize}
            \item Solusi optimal masalah LP dijamin berada pada titik ekstrem (*corner point*) ruang feasibel.
            \item Metode grafis sangat efektif untuk memahami geometri 2D sebelum beralih ke Algoritma Simpleks dimensi tinggi.
        \end{itemize}
    \end{tcolorbox}
\end{frame}

\end{document}
